\documentclass[11pt,a4paper]{article}
\usepackage[utf8]{inputenc}
\usepackage{amsmath,amssymb}
\usepackage{geometry}
\geometry{margin=2.5cm}

\title{Thermodynamic Efficiency in Distributed Computing}
\author{A. M. Wright \and E. R. Vance}
\date{October 2026}

\begin{document}
\maketitle

\begin{abstract}
We examine the fundamental thermodynamic bounds governing distributed state consensus protocols. By combining stochastic Langevin dynamics with informational entropy production rates, we derive minimum energy dissipation limits for arbitrary networked topologies.
\end{abstract}

\section{Introduction}
The fundamental relationship between physical entropy and logical information was pioneered by Shannon~\cite{shannon1948} and Landauer~\cite{landauer1961}. In decentralized computing, information reconciliation requires non-zero entropy production. Early investigations by Turing~\cite{turing1936} established computational bounds that motivate our current framework.

\section{Theoretical Formulation}
Consider a network of $N$ nodes connected via graph Laplacian $L$. The continuous state vector $\mathbf{x}(t)$ evolves according to:
\begin{equation}
  \frac{d\mathbf{x}}{dt} = -L \mathbf{x}(t) + \boldsymbol{\xi}(t)
  \label{eq:langevin}
\end{equation}
where $\boldsymbol{\xi}(t)$ is Gaussian white noise satisfying $\langle \boldsymbol{\xi}(t) \boldsymbol{\xi}^T(t') \rangle = 2 D I_N \delta(t - t')$.

The total entropy production rate evaluates to:
\begin{equation}
  \dot{S}_{\text{prod}} = \frac{\gamma}{T} \langle \dot{\mathbf{x}}^T \dot{\mathbf{x}} \rangle \ge 0
  \label{eq:entropy_rate}
\end{equation}

\section{Conclusion}
Our results confirm that respecting thermodynamic limits leads directly to more energy-efficient distributed consensus architectures.

\bibliographystyle{plain}
\bibliography{references}

\end{document}
